\chapter{Ruang Euklides}\label{ch:b1:euclid}

Menambahkan \omterm{def:b1:euclid:def}{hasil kali dalam} pada \omterm{def:b1:vspaces:def}{ruang vektor} real membeli gagasan
geometrinya --- panjang, sudut, keortogonalan, jarak --- dan satu
teorema yang menjulang di atas bab ini: bahwa setiap \omterm{def:b1:vspaces:subspace}{subruang} mempunyai
\omterm{thm:b1:euclid:projection}{proyeksi ortogonal}, yang dapat dihitung lewat Gram--Schmidt, dan mewujudkan
jarak terpendeknya. Adapun \omterm{def:b1:euclid:isometry}{isometri} bidang menutup bab ini dan
geometri tahun ini.

Sepanjang bab ini, $E$ merupakan \omterm{def:b1:vspaces:def}{ruang vektor} \emph{real}.

\section{Hasil kali dalam}

\begin{definition}\label{def:b1:euclid:def}
Sebuah \emph{hasil kali dalam}\index{hasil kali dalam} pada $E$ adalah \omterm{def:b1:logic:map}{pemetaan}
$\langle\cdot,\cdot\rangle \colon E \times E \to \R$ yang
bilinear, setangkup, dan tegas positif ($\langle x, x\rangle >
0$ untuk $x \neq 0$). Adapun ruang \omterm{def:b1:findim:def}{berdimensi hingga} yang dilengkapi demikian disebut
\emph{ruang Euklides}\index{ruang Euklides}. Sedangkan \emph{norma}
$x$ adalah $\norm{x} = \sqrt{\langle x, x\rangle}$, dan $d(x, y) =
\norm{x - y}$.
\end{definition}

\begin{example}\label{ex:b1:euclid:examples}
Pada $\R^n$: hasil kali kanoniknya $\langle x, y\rangle = \sum x_i
y_i$. Pada $C(\intcc{a}{b})$: $\langle f, g \rangle = \int_a^b fg$
(dengan ketegasan positifnya adalah \cref{thm:b1:integration:props}~(4)). Pada
$\R_n[X]$: $\langle P, Q\rangle = \int_0^1 PQ$, atau $\sum_{i}
P(x_i)Q(x_i)$ atas $n+1$ titik yang berbeda.
\end{example}

\begin{example}[Sudut antara dua polinomial]
\label{ex:b1:euclid:polyangle}
Sekali sebuah \omterm{def:b1:euclid:def}{hasil kali dalam} dipilih, \emph{sebarang} dua vektor taknol
mempunyai sudut, lewat $\cos\theta = \frac{\langle x,
y\rangle}{\norm x\,\norm y}$ (yaitu kosinus yang sah menurut
Cauchy--Schwarz). Bagi $X$ dan $X^2$ pada $\int_0^1$:
\[
\langle X, X^2\rangle = \frac14, \qquad
\norm X = \frac1{\sqrt3}, \qquad \norm{X^2} = \frac1{\sqrt5},
\qquad
\cos\theta = \frac{1/4}{1/\sqrt{15}} = \frac{\sqrt{15}}{4}
\approx 0.968 :
\]
yaitu sudut sekitar $14.5$ derajat --- jadi pada $\intcc{0}{1}$,
grafik $x$ dan $x^2$ ``nyaris sejajar'' dalam makna
rerata kuadratnya, dan itulah sebabnya membuang arah yang
terbagi itu (lewat Gram--Schmidt, di bawah) hanya menyisakan
koreksi kecil $X^2 - X + \frac16$.
\end{example}

\begin{theorem}[Cauchy--Schwarz; sifat normanya]\label{thm:b1:euclid:cs}
Untuk semua $x, y \in E$:\index{ketaksamaan Cauchy--Schwarz}
\[
\abs{\langle x, y\rangle} \leq \norm x\, \norm y ,
\]
dengan kesamaannya jika dan hanya jika $x, y$ sebanding. Akibatnya $\norm\cdot$
memenuhi ketaksamaan segitiga $\norm{x + y} \leq \norm x + \norm
y$ (dan $\norm{\lambda x} = \abs\lambda \norm x$, $\norm x = 0 \iff
x = 0$). Lebih-lebih:
\[
\norm{x+y}^2 = \norm x^2 + 2\langle x, y\rangle + \norm y^2,
\qquad
\langle x, y \rangle = \tfrac14\bigl(\norm{x+y}^2 -
\norm{x-y}^2\bigr) .
\]
\end{theorem}

\begin{proof}
Jika $y = 0$, segalanya sepele. Kalau tidak, kuadratik $t
\mapsto \norm{x + ty}^2 = \norm x^2 + 2t\langle x, y\rangle + t^2
\norm y^2$ bernilai $\geq 0$ untuk semua $t$: sehingga diskriminannya $\leq 0$,
dan itulah Cauchy--Schwarz; adapun kesamaannya berarti akar rangkap $t_0$, yaitu
$x + t_0 y = 0$ (lewat ketegasannya): jadi kesebandingan. Adapun ketaksamaan
segitiganya: jabarkanlah,
\[
\norm{x + y}^2 = \norm x^2 + 2\langle x, y\rangle + \norm y^2
\leq \norm x^2 + 2\norm x\,\norm y + \norm y^2
= \bigl(\norm x + \norm y\bigr)^2 ,
\]
dengan langkah tengahnya adalah Cauchy--Schwarz; dan kesamaannya memaksa $\langle
x, y\rangle = \norm x\norm y$, yaitu kasus kesamaan yang
\emph{positif}, jadi kesebandingan dengan nisbah taknegatif ---
sehingga secara geometri, segitiganya merosot hanya ketika kedua
vektornya menunjuk arah yang sama. Kedua kesamaan terakhirnya
merupakan penjabaran langsung (dengan yang kedua, yaitu
\emph{kesamaan polarisasi}, memulihkan hasil kalinya dari normanya).
\end{proof}

\section{Keortogonalan}

\begin{definition}\label{def:b1:euclid:orthogonal}
$x \perp y$ ketika $\langle x, y \rangle = 0$. Sebuah keluarga disebut
\emph{ortogonal} bila vektornya ortogonal berpasangan, dan
\emph{ortonormal}\index{keluarga ortonormal} bila lebih-lebih masing-masingnya
bernorma $1$. Adapun \emph{\omterm{def:b1:vspaces:sum}{pelengkap} ortogonal} sebuah \omterm{def:b1:vspaces:subspace}{subruang} $F$ adalah
\[
F^{\perp} = \{x \in E : \forall y \in F,\ \langle x, y\rangle =
0\},
\]
sebuah \omterm{def:b1:vspaces:subspace}{subruang} $E$.\index{pelengkap ortogonal}
\end{definition}

\begin{proposition}\label{prop:b1:euclid:orthfree}
(Pythagoras\index{teorema Pythagoras}) Jika $x \perp y$ maka
$\norm{x+y}^2 = \norm x^2 + \norm y^2$. Sebuah keluarga \omterm{def:b1:euclid:orthogonal}{ortogonal} berisi
vektor taknol bersifat \omterm{def:b1:vspaces:free}{bebas}. Pada \omterm{def:b1:vspaces:free}{basis} \omterm{def:b1:euclid:orthogonal}{ortonormal} $(e_1, \dots,
e_n)$, \omterm{prop:b1:vspaces:coordinates}{koordinat} dan hasil kalinya adalah
\[
x = \sum_{i} \langle x, e_i\rangle\, e_i,
\qquad
\langle x, y \rangle = \sum_i \langle x, e_i\rangle \langle y,
e_i\rangle,
\qquad
\norm x^2 = \sum_i \langle x, e_i\rangle^2 .
\]
\end{proposition}

\begin{proof}
Pythagoras: jabarkanlah. Kebebasannya: ambillah $\langle\,\cdot\,, x_j\rangle$
atas sebuah kombinasi nol: maka $\lambda_j \norm{x_j}^2 = 0$. \omterm{prop:b1:vspaces:coordinates}{Koordinatnya}:
tulislah $x = \sum \lambda_i e_i$ lalu ambillah hasil kalinya dengan $e_j$:
$\lambda_j = \langle x, e_j\rangle$; dan kedua rumusnya menyusul lewat
kebilinearan.
\end{proof}

\begin{example}[Koordinat ortonormal, dengan periksa Parseval]
\label{ex:b1:euclid:parsevalcheck}
Jabarkanlah $x = (1, 2, 3)$ pada \omterm{def:b1:vspaces:free}{basis} \omterm{def:b1:euclid:orthogonal}{ortonormal}
\cref{exo:b1:euclid:3},
\[
e_1 = \tfrac{1}{\sqrt2}(1,1,0), \quad
e_2 = \tfrac{1}{\sqrt6}(1,-1,2), \quad
e_3 = \tfrac{1}{\sqrt3}(-1,1,1).
\]
Tak ada sistem yang perlu dipecahkan --- hanya tiga \omterm{def:b1:euclid:def}{hasil kali dalam}:
\[
\langle x, e_1\rangle = \frac{3}{\sqrt2},
\qquad
\langle x, e_2\rangle = \frac{1 - 2 + 6}{\sqrt6} =
\frac{5}{\sqrt6},
\qquad
\langle x, e_3\rangle = \frac{-1 + 2 + 3}{\sqrt3} =
\frac{4}{\sqrt3}.
\]
Pengesahannya lewat rumus norma pada proposisinya:
\[
\frac{9}{2} + \frac{25}{6} + \frac{16}{3}
= \frac{27 + 25 + 32}{6} = 14 = \norm{x}^2 = 1 + 4 + 9 .
\]
Periksa jumlah kuadrat \omterm{prop:b1:vspaces:coordinates}{koordinat} ini (yaitu kesamaan Parseval yang
hingga) berongkos beberapa detik dan menangkap galat tanda dan
penormalan nyaris pasti --- jadikanlah ia kebiasaan kapan pun sebuah
penjabaran \omterm{def:b1:euclid:orthogonal}{ortonormal} dihitung; adapun versi berdimensi takhingganya,
bagi koefisien Fourier pada
\cref{ex:b1:euclid:trigortho}, merupakan teorema jilid Tahun~3.
\end{example}

\begin{theorem}[Gram--Schmidt]\label{thm:b1:euclid:gramschmidt}
Setiap \omterm{def:b1:euclid:def}{ruang Euklides} mempunyai \omterm{def:b1:vspaces:free}{basis} \omterm{def:b1:euclid:orthogonal}{ortonormal}. Secara eksplisit, dari sebarang
\omterm{def:b1:vspaces:free}{basis} $(v_1, \dots, v_n)$, resepnya\index{proses Gram--Schmidt}
\[
w_k = v_k - \sum_{i=1}^{k-1} \langle v_k, e_i\rangle\, e_i ,
\qquad
e_k = \frac{w_k}{\norm{w_k}}
\]
menghasilkan \omterm{def:b1:vspaces:free}{basis} \omterm{def:b1:euclid:orthogonal}{ortonormal} $(e_1, \dots, e_n)$ dengan
$\operatorname{Vect}(e_1, \dots, e_k) = \operatorname{Vect}(v_1,
\dots, v_k)$ untuk setiap $k$.
\end{theorem}

\begin{proof}
Induksi pada $k$. Dengan mengandaikan $(e_1, \dots, e_{k-1})$ \omterm{def:b1:euclid:orthogonal}{ortonormal}
dan merentang $\operatorname{Vect}(v_1, \dots, v_{k-1})$: maka vektor
$w_k$ \omterm{def:b1:euclid:orthogonal}{ortogonal} terhadap setiap $e_j$ ($j < k$) menurut konstruksinya
($\langle w_k, e_j \rangle = \langle v_k, e_j\rangle - \langle v_k,
e_j\rangle$), dan $w_k \neq 0$ karena $v_k \notin
\operatorname{Vect}(v_1, \dots, v_{k-1})$. Menormalkannya menjaga
keortogonalannya; sedangkan \omterm{def:b1:logic:statement}{pernyataan} \omterm{def:b1:vspaces:span}{rentangnya} berlaku karena $e_k$ merupakan
kombinasi atas $v_k$ dan $e_i$ yang terdahulu, secara terbalikkan.
\end{proof}

\begin{example}[Gram--Schmidt pada polinomial, selengkapnya]
\label{ex:b1:euclid:legendre}
Ortonormalkanlah $(1, X, X^2)$ di dalam $\R_2[X]$ dengan $\langle P,
Q\rangle = \int_0^1 PQ$. \emph{Langkah 1}: $\norm{1}^2 = 1$, sehingga
$e_1 = 1$. \emph{Langkah 2}: $w_2 = X - \langle X, 1\rangle\,1 = X -
\frac12$, dan $\norm{w_2}^2 = \int_0^1\bigl(x - \frac12\bigr)^2
\dd x = \frac1{12}$: jadi $e_2 = \sqrt{12}\,\bigl(X - \frac12\bigr)$.
\emph{Langkah 3}: $\langle X^2, e_1\rangle = \frac13$ dan
\[
\langle X^2, e_2\rangle
= \sqrt{12}\int_0^1 x^2\Bigl(x - \frac12\Bigr)\dd x
= \frac{\sqrt{12}}{12},
\qquad\text{sehingga}\qquad
w_3 = X^2 - \frac13 - \Bigl(X - \frac12\Bigr)
= X^2 - X + \frac16 .
\]
Normanya dihitung pada \cref{exo:b1:euclid:9}: $\norm{w_3}^2 =
\frac1{180}$, dan dari situlah $e_3 = \sqrt{180}\,\bigl(X^2 - X +
\frac16\bigr)$. Adapun \omterm{def:b1:poly:def}{polinomial} $1$, $X - \frac12$, $X^2 - X +
\frac16$ merupakan, sampai skalanya, \emph{\omterm{def:b1:poly:def}{polinomial} Legendre} yang pertama
pada \omterm{prop:b1:reals:intervals}{selang} $\intcc{0}{1}$; dan konstruksinya berlanjut satu
derajat setiap kali, dengan setiap \omterm{def:b1:poly:def}{polinomial} barunya \omterm{def:b1:euclid:orthogonal}{ortogonal} terhadap semua
pendahulunya. Perhatikanlah bagaimana algoritmanya mendaur ulang kerja terdahulunya:
\omterm{def:b1:linmaps:projection}{proyeksi} yang dikurangkan pada langkah 3 tepat merupakan hampiran afin
terbaik atas $X^2$ yang ditemukan pada \cref{ex:b1:euclid:bestapprox}
--- jadi Gram--Schmidt \emph{adalah} \omterm{thm:b1:euclid:projection}{proyeksi ortogonal} yang berulang.
\end{example}

\begin{theorem}[Proyeksi ortogonal]\label{thm:b1:euclid:projection}
Misalkan $F$ \omterm{def:b1:vspaces:subspace}{subruang} \omterm{def:b1:euclid:def}{ruang Euklides} $E$. Maka
\[
E = F \oplus F^{\perp},
\]
dan \omterm{def:b1:linmaps:projection}{proyeksi} $p_F$ yang berkaitan ke $F$ (yaitu \emph{proyeksi
ortogonal}\index{proyeksi ortogonal}) diberikan, pada sebarang
\omterm{def:b1:vspaces:free}{basis} \omterm{def:b1:euclid:orthogonal}{ortonormal} $(e_1, \dots, e_k)$ pada $F$, oleh $p_F(x) = \sum_i
\langle x, e_i\rangle e_i$. Ia mewujudkan jarak ke $F$:
karena untuk semua $y \in F$,
\[
\norm{x - p_F(x)} \leq \norm{x - y},
\]
dengan kesamaannya hanya untuk $y = p_F(x)$; dan orang menulis $d(x, F) = \norm{x -
p_F(x)}$.
\end{theorem}

\begin{proof}
Ambillah \omterm{def:b1:vspaces:free}{basis} \omterm{def:b1:euclid:orthogonal}{ortonormal} $(e_i)_{i \leq k}$ pada $F$
(\cref{thm:b1:euclid:gramschmidt} di dalam $F$) lalu tetapkanlah $\pi(x) =
\sum \langle x, e_i\rangle e_i \in F$. Maka $x - \pi(x) \perp e_j$
untuk setiap $j$ (lewat penghapusan yang sama seperti di atas), sehingga $x - \pi(x) \in
F^\perp$: jadi $E = F + F^\perp$. Dan $F \cap F^\perp = \{0\}$: karena
vektor yang demikian memenuhi $\langle x, x\rangle = 0$. Jadi jumlahnya langsung dan
$\pi = p_F$.

Jaraknya: untuk $y \in F$, uraikanlah $x - y = (x - p_F(x)) + (p_F(x) -
y)$, yaitu potongan \omterm{def:b1:euclid:orthogonal}{ortogonal} ($F^\perp$ dan $F$); lalu Pythagoras:
\[
\norm{x - y}^2 = \norm{x - p_F(x)}^2 + \norm{p_F(x) - y}^2
\geq \norm{x - p_F(x)}^2,
\]
dengan kesamaannya jika dan hanya jika $y = p_F(x)$.
\end{proof}

\begin{example}[Proyeksi tak pernah memanjangkan]
\label{ex:b1:euclid:bessel}
Menerapkan Pythagoras pada pembelahan $x = p_F(x) + (x - p_F(x))$:
\[
\norm{p_F(x)}^2 = \norm x^2 - \norm{x - p_F(x)}^2 \leq
\norm x^2 ,
\]
dengan kesamaannya jika dan hanya jika $x \in F$. Pada \omterm{def:b1:vspaces:free}{basis} \omterm{def:b1:euclid:orthogonal}{ortonormal} $(e_1,
\dots, e_k)$ pada $F$ ini terbaca $\sum_{i \leq k}\langle x,
e_i\rangle^2 \leq \norm x^2$ (yaitu \emph{ketaksamaan Bessel}):
sebanyak apa pun arah \omterm{def:b1:euclid:orthogonal}{ortonormal} yang diukur orang, kuadrat
\omterm{prop:b1:vspaces:coordinates}{koordinatnya} tak pernah melampaui kuadrat panjangnya --- bandingkanlah
kesamaan persis pada \cref{ex:b1:euclid:parsevalcheck} ketika
keluarganya \omterm{def:b1:vspaces:free}{basis} penuh. Ketaksamaan sebaris inilah yang membuat
koefisien Fourier terjumlahkan pada jilid Tahun~3; sedangkan di sini ia
sudah menjelaskan mengapa menambahkan lebih banyak fungsi \omterm{def:b1:vspaces:free}{basis} pada suatu
pencocokan kuadrat terkecil hanya dapat mengurangi sisaannya.
\end{example}

\begin{example}[Hampiran kuadratik terbaik]\label{ex:b1:euclid:bestapprox}
Di dalam $C(\intcc{0}{1})$ dengan $\langle f, g\rangle = \int_0^1 fg$,
\omterm{def:b1:poly:def}{polinomial} berderajat $\leq 1$ yang terdekat dengan $f(x) = x^2$ dalam
jarak (rerata kuadrat) yang berkaitan adalah $p_F(f)$ dengan $F = \R_1[X]$.
Gram--Schmidt pada $(1, X)$: $e_1 = 1$, $w_2 = X - \frac12$,
$\norm{w_2}^2 = \int_0^1 (x - \frac12)^2 = \frac{1}{12}$, $e_2 =
\sqrt{12}\,(X - \tfrac12)$. Lalu
\[
p_F(f) = \langle f, e_1\rangle e_1 + \langle f, e_2\rangle e_2
= \frac13 + 12\Bigl(\int_0^1 x^2\bigl(x - \tfrac12\bigr)\dd
x\Bigr)\bigl(X - \tfrac12\bigr)
= X - \frac{1}{6},
\]
dengan memakai $\int_0^1 x^2(x - \frac12)\dd x = \frac14 - \frac16 =
\frac{1}{12}$. Jadi gagasan ``kuadrat terkecil'' dalam satu baris aljabar
\omterm{def:b1:linmaps:def}{linear}.
\end{example}

\begin{method}[Tiga rute menuju sebuah jarak $d(x, F)$]
\label{met:b1:euclid:distance}
\begin{enumerate}
  \item \emph{\omterm{def:b1:vspaces:free}{Basis} \omterm{def:b1:euclid:orthogonal}{ortonormal} pada $F$}: maka $p_F(x) = \sum_i
        \langle x, e_i\rangle e_i$ dan, menurut Pythagoras,
        \[
        d(x, F)^2 = \norm{x}^2 - \norm{p_F(x)}^2
        = \norm x^2 - \sum_i \langle x, e_i\rangle^2 ,
        \]
        yang kerap lebih murah daripada menghitung $x - p_F(x)$ itu sendiri.
  \item \emph{Persamaan normal}: dengan sebarang keluarga pembangun
        $F$, pecahkanlah $\langle x - p, v_j\rangle = 0$ bagi
        koefisien $p$ (\cref{exo:b1:euclid:5}) --- jadi tak ada
        pengortonormalan yang diperlukan.
  \item \emph{Lewat \omterm{def:b1:vspaces:sum}{pelengkapnya}}: jika $F^\perp$ lebih kecil
        daripada $F$ (misalnya $F$ sebuah \omterm{def:b1:linmaps:forms}{hiperbidang}, dan $F^\perp$ sebuah garis
        $\operatorname{Vect}(n)$), proyeksikanlah ke $F^\perp$
        sebagai gantinya:
        \[
        d(x, F) = \norm{p_{F^\perp}(x)} =
        \frac{\abs{\langle x, n\rangle}}{\norm n} ,
        \]
        yaitu rumus jarak ke sebuah bidang yang klasik
        (\cref{exo:b1:multivar:8} memakainya).
\end{enumerate}
Rute 3 merupakan kasus khusus sebuah gerak naluri umum: yaitu selalu proyeksikanlah
ke mana pun di antara $F$, $F^\perp$ yang berdimensi lebih kecil.
\end{method}

\begin{example}[Keortogonalan trigonometri: pratinjau Fourier]
\label{ex:b1:euclid:trigortho}
Pada $C(\intcc{0}{2\pi})$ dengan $\langle f, g\rangle =
\frac1\pi\int_0^{2\pi} fg$, keluarga
\[
\Bigl(\frac{1}{\sqrt2},\ \cos x,\ \sin x,\ \cos 2x,\ \sin 2x,\
\dots\Bigr)
\]
bersifat \omterm{def:b1:euclid:orthogonal}{ortonormal}: misalnya $\langle \cos px, \cos qx\rangle =
\frac1\pi\int_0^{2\pi}\cos px\cos qx\,\dd x = 0$ untuk $p \neq q$
(linearkanlah hasil kalinya menjadi $\frac12[\cos(p{-}q)x +
\cos(p{+}q)x]$ lalu integralkanlah atas periode penuhnya), sedangkan
$\frac1\pi\int_0^{2\pi}\cos^2 px\,\dd x = 1$. Maka \omterm{thm:b1:euclid:projection}{proyeksi
ortogonal} ke \omterm{def:b1:vspaces:span}{rentang} $2N + 1$ fungsi yang pertama ini
karenanya berkoordinat $\langle f, e_i\rangle$ --- yaitu \omterm{thm:b1:integration:def}{integral}
terhadap kosinus dan sinus. Inilah \emph{koefisien
Fourier} pada $f$, dan \omterm{def:b1:linmaps:projection}{proyeksinya} adalah hampiran
trigonometri rerata kuadrat terbaiknya; adapun jilid Tahun~3
mempelajari kekonvergenannya. Keortogonalanlah yang mengerjakan segalanya:
karena rumus koefisiennya adalah
\cref{thm:b1:euclid:projection} secara harfiah.
\end{example}

\section{Isometri bidang}

\begin{definition}\label{def:b1:euclid:isometry}
Sebuah endomorfisma $u$ pada \omterm{def:b1:euclid:def}{ruang Euklides} disebut
\emph{isometri}\index{isometri} (atau \omterm{def:b1:logic:map}{pemetaan} \omterm{def:b1:euclid:orthogonal}{ortogonal}) ketika ia
mengawetkan normanya: $\norm{u(x)} = \norm x$ untuk semua $x$ ---
setara dengan itu (lewat polarisasi) ia mengawetkan \omterm{def:b1:euclid:def}{hasil kali dalamnya};
setara dengan itu matriksnya $A$ pada \omterm{def:b1:vspaces:free}{basis} \omterm{def:b1:euclid:orthogonal}{ortonormal} memenuhi
$A^{\mathsf T} A = I$. Adapun isometri membentuk \omterm{def:b1:structures:group}{grup}, yaitu \omterm{def:b1:structures:group}{grup} \omterm{def:b1:euclid:orthogonal}{ortogonal}
$O(E)$.
\end{definition}

\begin{example}[Mengenali sebuah isometri sekali pandang]
\label{ex:b1:euclid:recognize}
Apakah $A = \dfrac15\begin{pmatrix} 3 & -4\\ 4 & 3\end{pmatrix}$
\omterm{def:b1:euclid:orthogonal}{ortogonal}? Kolomnya: normanya $\frac15\sqrt{9 + 16} = 1$ dan
$\frac15\sqrt{16 + 9} = 1$; hasil kalinya $\frac1{25}(3\cdot(-4) +
4\cdot3) = 0$. Ya --- dan $\det A = \frac{9 + 16}{25} = 1$, sehingga
ia rotasi $R_\theta$ dengan $\cos\theta = \frac35$,
$\sin\theta = \frac45$ (yaitu ``rotasi $3$-$4$-$5$'', yang
sudutnya bukan pecahan $\pi$ yang istimewa). Sebaliknya $B =
\frac{1}{\sqrt2}\begin{pmatrix} 1 & 1\\ 0 & 1\end{pmatrix}$
tampak berdeterminan terskalakan satuan tetapi kolom pertamanya tak bersatuan
($\frac1{\sqrt2}$): jadi tak \omterm{def:b1:euclid:orthogonal}{ortogonal} --- karena determinan $\pm1$ belaka
tak mengesahkan apa pun, kolomnya mesti diperiksa.
\end{example}

\begin{theorem}[Isometri bidang]\label{thm:b1:euclid:plane}
Pada \omterm{def:b1:vspaces:free}{basis} \omterm{def:b1:euclid:orthogonal}{ortonormal} sebuah bidang Euklides, matriks
\omterm{def:b1:euclid:isometry}{isometrinya} tepat berupa
\[
R_\theta = \begin{pmatrix}
\cos\theta & -\sin\theta\\ \sin\theta & \cos\theta
\end{pmatrix}
\quad (\text{rotasi bersudut } \theta,\ \det = 1),
\]
\[
S_\theta = \begin{pmatrix}
\cos\theta & \sin\theta\\ \sin\theta & -\cos\theta
\end{pmatrix}
\quad (\det = -1),
\]
dengan yang terakhirnya adalah pencerminan pada garis yang bersudut
$\frac\theta2$ dengan vektor \omterm{def:b1:vspaces:free}{basis} pertamanya.
\end{theorem}

\begin{proof}
Misalkan $A = \begin{pmatrix} a & c\\ b & d\end{pmatrix}$ dengan
$A^{\mathsf T}A = I$: maka kolomnya bersatuan dan \omterm{def:b1:euclid:orthogonal}{ortogonal}. Kolom
pertamanya adalah $(\cos\theta, \sin\theta)$ bagi suatu $\theta$; sedangkan yang kedua,
yang bersatuan dan \omterm{def:b1:euclid:orthogonal}{ortogonal} terhadapnya, adalah $\pm(-\sin\theta, \cos\theta)$. Adapun
tanda $+$ memberikan $R_\theta$; sedangkan tanda $-$ memberikan $S_\theta$. Orang memeriksa
$S_\theta^2 = I$ dan bahwa vektor $(\cos\frac\theta2,
\sin\frac\theta2)$ tertetapkan sedangkan \omterm{def:b1:euclid:orthogonal}{ortogonalnya} terbalikkan: jadi sebuah
pencerminan. (Dan $R_\alpha R_\beta = R_{\alpha+\beta}$: sehingga \omterm{def:b1:structures:group}{grup}
rotasinya adalah \omterm{def:b1:structures:group}{grup} sudutnya --- bandingkanlah
\cref{thm:b1:complex:funceq}.)
\end{proof}

\begin{omfigure}
\begin{tikzpicture}[scale=1.15]
  \draw[->, thin] (-1.2,0) -- (2.9,0) node[below] {$x$};
  \draw[->, thin] (0,-1.2) -- (0,2.6) node[left] {$y$};
  \draw[omDef, very thick] (-0.9,0) -- (2.6,0)
    node[below, pos=0.9, font=\small] {sumbu $r_1$};
  \draw[omThm, very thick] (-0.9,-0.9) -- (2.3,2.3)
    node[right, pos=0.97, font=\small] {sumbu $r_2$};
  \fill (2,0.5) circle (0.05) node[right, font=\small] {$M$};
  \fill (2,-0.5) circle (0.05)
    node[right, font=\small] {$r_1(M)$};
  \fill (-0.5,2) circle (0.05)
    node[above, font=\small] {$r_2(r_1(M))$};
  \draw[omDef, dashed] (2,0.5) -- (2,-0.5);
  \draw[omThm, dashed] (2,-0.5) -- (-0.5,2);
  \draw[omProp, thick, ->] (0.9,0) arc (0:45:0.9);
  \node[omProp, font=\small] at (1.15,0.32) {$\frac\pi4$};
\end{tikzpicture}

{\small Dua pencerminan membuat sebuah rotasi: mencerminkan $M = (2,
0.5)$ pada sumbu $x$, lalu pada garis $y = x$, mendarat di
$(-0.5, 2)$ --- yaitu peta $M$ di bawah rotasi bersudut
$\frac\pi2$ terhadap titik asalnya, yakni \emph{dua kali} sudut
$\frac\pi4$ antara sumbunya. Adapun soal akhir pekannya mengubah
gambar ini menjadi hukum komposisi semua \omterm{def:b1:euclid:isometry}{isometri} bidang.}
\end{omfigure}

\begin{remark}[Jebakan yang lazim]
\label{rem:b1:euclid:pitfalls}
\emph{Rumus \omterm{def:b1:linmaps:projection}{proyeksinya} menuntut \omterm{def:b1:vspaces:free}{basis} \omterm{def:b1:euclid:orthogonal}{ortonormal}}: karena bagi
keluarga pembangun $(v_i)$ pada $F$ belaka, jumlah $\sum_i\langle
x, v_i\rangle v_i$ \emph{bukan} $p_F(x)$ (ujilah $F = \R^2$, $v_1
= e_1$, $v_2 = e_1 + e_2$); jadi dengan keluarga yang tak \omterm{def:b1:euclid:orthogonal}{ortonormal}, pecahkanlah
persamaan normalnya sebagai gantinya
(\cref{met:b1:euclid:distance}~(2)).
\emph{Keluarga \omterm{def:b1:euclid:orthogonal}{ortogonal} mesti menghindari $0$ agar \omterm{def:b1:vspaces:free}{bebas}}: karena vektor
nol \omterm{def:b1:euclid:orthogonal}{ortogonal} terhadap segalanya, termasuk terhadap dirinya sendiri ---
sedangkan kebebasan pada \cref{prop:b1:euclid:orthfree} menuntut vektor
yang taknol.
\emph{$F^\perp$ bergantung pada \omterm{def:b1:euclid:def}{hasil kali dalamnya}}: di dalam $\R_1[X]$,
\omterm{def:b1:vspaces:sum}{pelengkap} $\operatorname{Vect}(X)$ bagi $\int_0^1 PQ$ bukanlah
konstantanya melainkan $\operatorname{Vect}(1 - \frac32 X)$ ---
hitunglah $\int_0^1 x(1 - \frac32 x) = \frac12 - \frac12 = 0$;
jadi ``tegak lurus'' tak bermakna sampai hasil kalinya dinamai.
\emph{Jangan menjabarkan $\norm{x + y}$ secara \omterm{def:b1:linmaps:def}{linear}}: karena kesamaan
yang benar adalah $\norm{x+y}^2 = \norm x^2 + 2\langle x, y\rangle +
\norm y^2$; sedangkan suku silangnya lenyap hanya di bawah keortogonalan
(Pythagoras), dan ketaksamaan segitiganya memang sebuah
\emph{ketaksamaan}.
\emph{Mengirim vektor satuan ke vektor satuan belumlah cukup}: karena $u(x,
y) = (x + y,\ 0)$ memetakan kedua vektor \omterm{def:b1:vspaces:free}{basis} kanoniknya ke vektor
satuan $(1, 0)$, namun $\norm{u(1,1)} = 2 \neq \sqrt2$: jadi bukan
\omterm{def:b1:euclid:isometry}{isometri}. Definisinya menuntut $\norm{u(x)} = \norm x$ untuk
\emph{semua} $x$; dan dalam istilah matriks $A^{\mathsf T}A = I$, yaitu
kolom yang bersatuan \emph{dan \omterm{def:b1:euclid:orthogonal}{ortogonal} berpasangan} --- jadi kedua
syaratnya, diperiksa bersama.
\end{remark}

\begin{remark}[Ke mana hasil kali dalamnya pergi]
\label{rem:b1:euclid:whereused}
\omterm{thm:b1:euclid:projection}{Proyeksi ortogonal} merupakan teorema bab ini yang paling banyak diterapkan:
ia melandasi kuadrat terkecil (karena soal akhir pekan
\cref{ch:b1:multivar} membangun garis regresi di atasnya), koefisien
Fourier (\cref{ex:b1:euclid:trigortho}), dan persamaan
normal \cref{exo:b1:euclid:5}, yang dipecahkan analisis numeris
secara besar-besaran. Adapun penggolongan $R_\theta / S_\theta$
dilengkapkan di bawah: karena soal akhir pekannya menggolongkan \emph{semua}
transformasi bidang yang mengawetkan jarak, yang \omterm{def:b1:linmaps:def}{linear} ataupun bukan,
beserta \omterm{def:b1:structures:group}{grup} hingganya --- yaitu matematika rozet dan
segi banyak beraturan. Pada jilid Tahun~2 \omterm{def:b1:euclid:def}{hasil kali dalamnya} bertemu
teori nilai eigen (yaitu matriks setangkup, bentuk kuadratik); sedangkan pada
Tahun~3, geometri Euklides berdimensi takhingga menjadi teori ruang
Hilbert.
\end{remark}

\begin{remark}[Cakrawala di dalam Buku 3]
\label{rem:b1:euclid:perspectives}
Dua jembatan meninggalkan bab ini. \emph{Ke belakang}, menuju aljabar
\omterm{def:b1:linmaps:def}{linearnya}: matriks Gram pada \cref{exo:b1:euclid:11} membungkus
\omterm{def:b1:euclid:def}{hasil kali dalam} ke dalam mesin determinan
\cref{ch:b1:det}, dan \omterm{thm:b1:euclid:projection}{proyeksi ortogonal} merupakan
proyektor istimewa \cref{ch:b1:linmaps} yang kernelnya
$F^\perp$ --- sehingga seluruh aljabarnya ($p^2 = p$, $s = 2p -
\mathrm{id}$) berlaku secara harfiah, kini dengan bonus bahwa $\norm{x
- p(x)}$ merupakan sebuah \emph{jarak}. \emph{Ke depan}, menuju analisisnya:
\cref{ch:b1:curves} mengukur panjang busur dengan norma bab ini
dan tak menggolongkan apa pun tanpa \omterm{def:b1:euclid:isometry}{isometrinya};
\cref{ch:b1:multivar} membaca gradiennya lewat
Cauchy--Schwarz (yaitu pendakian tercuram) dan menutup jilid ini dengan
kuadrat terkecil, yang merupakan \cref{thm:b1:euclid:projection} yang diterapkan
pada sebuah vektor data di dalam $\R^n$. Jadi \omterm{def:b1:euclid:def}{hasil kali dalamnya} adalah titik
tempat aljabar buku ini dan analisisnya akhirnya bertemu.
\end{remark}

\section{Latihan}

\begin{exercise}[$\star$]\label{exo:b1:euclid:1}
Di dalam $\R^3$ yang kanonik: hitunglah $\langle u, v\rangle$, $\norm u$,
$\norm v$ dan sudut antara $u = (1, 2, 2)$ dan $v = (2, -2,
1)$. Periksalah Cauchy--Schwarz secara numeris.
\end{exercise}

\begin{exercise}[$\star$]\label{exo:b1:euclid:2}
Buktikan kesamaan jajargenjang $\norm{x+y}^2 + \norm{x-y}^2 =
2\norm x^2 + 2\norm y^2$ pada sebarang \omterm{def:b1:euclid:def}{ruang Euklides}, lalu pakailah ia untuk menunjukkan
bahwa norma sup pada $\R^2$, $\norm{(x,y)}_\infty = \max(\abs x,
\abs y)$, tak berasal dari sebuah \omterm{def:b1:euclid:def}{hasil kali dalam}.
\end{exercise}

\begin{exercise}[$\star$]\label{exo:b1:euclid:3}
Terapkanlah Gram--Schmidt pada $\bigl((1,1,0), (1,0,1), (0,1,1)\bigr)$ di dalam
$\R^3$ yang kanonik.
\end{exercise}

\begin{exercise}[$\star$]\label{exo:b1:euclid:4}
Di dalam $\R^3$, misalkan $F = \operatorname{Vect}\bigl((1,1,1)\bigr)$.
Tentukanlah $F^\perp$ (persamaan dan basisnya), matriks $p_F$ pada
\omterm{def:b1:vspaces:free}{basis} kanoniknya, dan $d\bigl((1, 2, 3), F\bigr)$.
\end{exercise}

\begin{exercise}[$\star\star$]\label{exo:b1:euclid:5}
(Persamaan normal) Misalkan $F =
\operatorname{Vect}\bigl((1,0,1),(0,1,1)\bigr) \subseteq \R^3$ dan
$x = (1, 1, 4)$. Hitunglah $p_F(x)$ dengan memecahkan $\langle x - p, v
\rangle = 0$ bagi kedua pembangunnya ($p = \alpha(1,0,1) +
\beta(0,1,1)$), lalu $d(x, F)$. Mengapakah Gram--Schmidt tak diperlukan
di sini?
\end{exercise}

\begin{exercise}[$\star\star$]\label{exo:b1:euclid:6}
Untuk $f$ yang \omterm{def:b1:continuity:continuous}{kontinu} pada $\intcc{0}{1}$, buktikanlah
\[
\Bigl(\int_0^1 f\Bigr)^{\!2} \leq \int_0^1 f^2 ,
\]
beserta kasus kesamaannya, sebagai contoh Cauchy--Schwarz di dalam
$C(\intcc{0}{1})$. Lalu buktikanlah $\bigl(\sum_{i=1}^n a_i\bigr)^2 \leq
n \sum a_i^2$ bagi $a_i$ yang real.
\end{exercise}

\begin{exercise}[$\star\star$]\label{exo:b1:euclid:7}
Buktikan bahwa untuk setiap \omterm{def:b1:vspaces:subspace}{subruang} $F$ pada sebuah \omterm{def:b1:euclid:def}{ruang Euklides}:
$(F^{\perp})^{\perp} = F$ dan $\dim F^{\perp} = \dim E - \dim F$.
\end{exercise}

\begin{exercise}[$\star\star$]\label{exo:b1:euclid:8}
Kenalilah \omterm{def:b1:euclid:isometry}{isometri} bidang yang bermatriks
\[
A = \frac{1}{\sqrt 2}\begin{pmatrix} 1 & -1\\ 1 & 1\end{pmatrix},
\qquad
B = \frac{1}{5}\begin{pmatrix} 3 & 4\\ 4 & -3\end{pmatrix}
\]
(jenis, sudut atau sumbunya). Hitunglah $A^8$ dan $B^2$ tanpa
mengalikan matriks.
\end{exercise}

\begin{exercise}[$\star\star\star$]\label{exo:b1:euclid:9}
(Minimum sebagai sebuah \omterm{def:b1:linmaps:projection}{proyeksi}) Hitunglah
\[
\min_{(a, b) \in \R^2} \int_0^1 \bigl(x^2 - a - bx\bigr)^2 \dd x ,
\]
dengan memakai \cref{ex:b1:euclid:bestapprox}: karena minimumnya adalah $\norm{f -
p_F(f)}^2$ bagi $f = X^2$, $F = \R_1[X]$.
\end{exercise}

\begin{exercise}[$\star\star\star$]\label{exo:b1:euclid:10}
Misalkan $u$ sebuah \omterm{def:b1:euclid:isometry}{isometri} pada \omterm{def:b1:euclid:def}{ruang Euklides} $E$. Buktikan bahwa
$\ker(u - \mathrm{id}) \perp \operatorname{im}(u - \mathrm{id})$,
lalu simpulkanlah $E = \ker(u - \mathrm{id}) \oplus \operatorname{im}(u -
\mathrm{id})$. \emph{(Hitunglah $\langle x - u(x), y\rangle$ bagi
$u(y) = y$, dengan memakai pengawetan hasil kalinya.)}
\end{exercise}

\begin{exercise}[$\star\star$]\label{exo:b1:euclid:11}
(Matriks Gram) Untuk vektor $v_1, \dots, v_k$ pada sebuah \omterm{def:b1:euclid:def}{ruang
Euklides}, misalkan $G = \bigl(\langle v_i, v_j\rangle\bigr)_{1 \leq i, j
\leq k}$ \emph{matriks Gram}-nya.
\begin{enumerate}
  \item Buktikan bahwa $(v_1, \dots, v_k)$ \omterm{def:b1:vspaces:free}{bebas} jika dan hanya jika
        $G$ dapat dibalik. \emph{(Jika $Gc = 0$, hitunglah
        $\norm{\sum_i c_i v_i}^2$.)}
  \item Hitunglah matriks Gram atas $(1, X, X^2)$ di dalam $\R_2[X]$
        dengan $\langle P, Q\rangle = \int_0^1 PQ$, kenalilah
        \omterm{pb:b1:det:1}{matriks Hilbert} $H_3$ pada soal akhir pekan
        \cref{ch:b1:det}, lalu simpulkanlah kebebasannya dari $\det H_3 =
        \frac1{2160} \neq 0$.
\end{enumerate}
\end{exercise}

\begin{exercise}[$\star\star\star$]\label{exo:b1:euclid:12}
Misalkan $u$ sebuah endomorfisma pada \omterm{def:b1:euclid:def}{ruang Euklides} $E$ yang matriksnya
$A$ pada \omterm{def:b1:vspaces:free}{basis} \omterm{def:b1:euclid:orthogonal}{ortonormal} sekaligus \omterm{def:b1:euclid:orthogonal}{ortogonal} ($A^{\mathsf T}A
= I$) dan setangkup ($A^{\mathsf T} = A$).
\begin{enumerate}
  \item Tunjukkan $A^2 = I$, lalu simpulkanlah (lewat
        \cref{thm:b1:linmaps:projchar}) bahwa $E = \ker(u -
        \mathrm{id}) \oplus \ker(u + \mathrm{id})$.
  \item Tunjukkan bahwa kedua \omterm{def:b1:vspaces:subspace}{subruangnya} \omterm{def:b1:euclid:orthogonal}{ortogonal}, sehingga $u$
        merupakan \emph{kesetangkupan \omterm{def:b1:euclid:orthogonal}{ortogonal}} terhadap $F =
        \ker(u - \mathrm{id})$: yaitu pencerminan lewat $F$.
        \emph{(Bagi $u(x) = x$ dan $u(y) = -y$, hitunglah $\langle
        x, y\rangle$ dengan dua cara.)}
  \item Golongkanlah kasus bidangnya: matriks yang mana pada
        \cref{thm:b1:euclid:plane} yang setangkup, dan apakah
        \omterm{def:b1:logic:map}{pemetaan} yang bersesuaian dengannya?
\end{enumerate}
\end{exercise}

\section{Soal: isometri bidang dan teorema Leonardo}

\begin{problem}\label{pb:b1:euclid:1}
Sebuah \emph{\omterm{def:b1:euclid:isometry}{isometri} bidang} adalah sebarang \omterm{def:b1:logic:map}{pemetaan} $f \colon \R^2 \to
\R^2$ yang mengawetkan jarak: $\norm{f(x) - f(y)} = \norm{x - y}$
untuk semua $x, y$ --- tanpa mengandaikan kelinearan. Soal ini membuktikan
bahwa \omterm{def:b1:logic:map}{pemetaan} yang demikian tepat berupa translasi, rotasi,
pencerminan dan \omterm{pb:b1:euclid:1}{pencerminan geser} (yaitu penggolongan \omterm{def:b1:euclid:isometry}{isometri}
bidang), menghitung komposisinya, lalu menentukan semua
\omterm{def:b1:structures:group}{grup} \emph{hingga}-nya: yaitu teorema Leonardo, matematika
di balik pola rozet. Sejak Bagian II kita menyamakan $\R^2$ dengan
$\C$ (\cref{ch:b1:complex}): dengan \omterm{def:b1:euclid:def}{hasil kali dalam} kanoniknya adalah
$\langle z, w\rangle = \operatorname{Re}(z\conj w)$ dan normanya
adalah modulusnya.\index{isometri}\index{pencerminan geser}
\index{teorema Leonardo}\index{grup dihedral}

\medskip
\noindent\textbf{Bagian I --- Setiap \omterm{def:b1:euclid:isometry}{isometri} bersifat afin.}
\begin{enumerate}
  \item Periksalah bahwa translasi $t_a(x) = x + a$, \omterm{def:b1:euclid:isometry}{isometri}
        \omterm{def:b1:linmaps:def}{linear}, dan semua komposisinya merupakan \omterm{def:b1:euclid:isometry}{isometri},
        dan bahwa \omterm{def:b1:euclid:isometry}{isometrinya} membentuk \omterm{def:b1:structures:group}{grup} di bawah komposisi.
  \item Misalkan $f$ sebuah \omterm{def:b1:euclid:isometry}{isometri} dengan $f(0) = 0$. Tunjukkan bahwa $f$
        mengawetkan norma, lalu --- lewat polarisasi,
        \cref{thm:b1:euclid:cs} --- bahwa $\langle f(x),
        f(y)\rangle = \langle x, y\rangle$ untuk semua $x, y$.
  \item Masih dengan $f(0) = 0$: jabarkanlah $\norm{f(x + y) - f(x) -
        f(y)}^2$ dan $\norm{f(\lambda x) - \lambda f(x)}^2$
        dengan memakai pertanyaan 2, lalu simpulkanlah bahwa $f$ bersifat
        \emph{\omterm{def:b1:linmaps:def}{linear}}: yaitu $f \in O(\R^2)$.
  \item Simpulkanlah bahwa setiap \omterm{def:b1:euclid:isometry}{isometri} $f$ tertulis \emph{secara tunggal} sebagai
        $f = t_a \circ g$ dengan $a = f(0)$ dan $g$ sebuah \omterm{def:b1:euclid:isometry}{isometri}
        \omterm{def:b1:linmaps:def}{linear} (yaitu \emph{bagian \omterm{def:b1:linmaps:def}{linear}} pada $f$).
  \item Sebutlah $f$ \emph{langsung} bila $\det g = 1$, dan \emph{taklangsung}
        bila $\det g = -1$. Tunjukkan bahwa bagian \omterm{def:b1:linmaps:def}{linear} sebuah
        komposisi adalah komposisi bagian \omterm{def:b1:linmaps:def}{linearnya}, lalu
        nyatakanlah kaidah tanda yang dihasilkannya
        (yaitu langsung/taklangsung berkomposisi bagai $+1/-1$).
\end{enumerate}

\medskip
\noindent\textbf{Bagian II --- Keempat jenisnya.} Lewat
\cref{thm:b1:euclid:plane}, \omterm{def:b1:euclid:isometry}{isometri} \omterm{def:b1:linmaps:def}{linear} pada $\C$ adalah $z
\mapsto az$ dan $z \mapsto a\conj z$ dengan $\abs a = 1$; sehingga setiap
\omterm{def:b1:euclid:isometry}{isometri} bidang berupa
\[
f(z) = a z + b \quad (\text{langsung}) \qquad\text{atau}\qquad
f(z) = a\conj z + b \quad (\text{taklangsung}), \qquad \abs a = 1 .
\]
\begin{enumerate}[resume]
  \item Periksalah kamusnya: bahwa $R_\theta$ adalah $z \mapsto
        \eu^{\iu\theta}z$ dan $S_\theta$ adalah $z \mapsto
        \eu^{\iu\theta}\conj z$ (periksalah keduanya pada $1$ dan $\iu$).
  \item (Kasus langsungnya) Misalkan $f(z) = az + b$, $\abs a = 1$. Tunjukkan:
        bila $a = 1$, maka $f$ sebuah translasi; bila $a \neq 1$, maka $f$ mempunyai
        titik tetap tunggal $z_0 = b/(1 - a)$ dan $f(z) - z_0
        = a(z - z_0)$: yaitu \emph{rotasi} berpusat $z_0$ dan
        bersudut $\arg a$.
  \item (Kasus taklangsungnya) Misalkan $f(z) = a\conj z + b$, dan $v =
        a\conj b + b$. Tunjukkan $f \circ f = t_v$ dan $f \circ t_v =
        t_v \circ f$. Jika $v = 0$: tunjukkanlah bahwa titik tengah $z$
        dan $f(z)$ merupakan titik tetap, dan bahwa $f$ sebuah
        \emph{pencerminan} pada sebuah garis. Jika $v \neq 0$: tunjukkanlah bahwa $r
        = t_{-v/2}\circ f$ sebuah pencerminan yang sumbunya
        sejajar $v$, sehingga $f = t_{v/2} \circ r$ merupakan
        \emph{\omterm{pb:b1:euclid:1}{pencerminan geser}}. Simpulkanlah: bahwa setiap \omterm{def:b1:euclid:isometry}{isometri} bidang
        berupa translasi, rotasi, pencerminan atau \omterm{pb:b1:euclid:1}{pencerminan
        geser} (yaitu penggolongan \omterm{def:b1:euclid:isometry}{isometri} bidang).
  \item (Komposisinya) Tunjukkan: bahwa komposisi rotasi
        bersudut $\alpha$ dan $\beta$ adalah rotasi bersudut
        $\alpha + \beta$ (yaitu translasi bila $\alpha + \beta \in
        2\pi\Z$); sedangkan komposisi dua pencerminan adalah
        rotasi bersudut dua kali sudut antara sumbunya (yaitu
        translasi bila sumbunya sejajar).
  \item Simpulkanlah bahwa setiap \omterm{def:b1:euclid:isometry}{isometri} bidang merupakan komposisi paling
        banyak tiga pencerminan.
\end{enumerate}

\medskip
\noindent\textbf{Bagian III --- Tiga pengenalan.}
\begin{enumerate}[resume]
  \item Golongkanlah selengkapnya $f(z) = \iu\conj z + 1 + \iu$:
        jenis, sumbu, dan vektor gesernya.
  \item Misalkan $f$ rotasi bersudut $\frac\pi2$ terhadap $0$
        dan $g$ rotasi bersudut $\frac\pi2$ terhadap $1$.
        Hitunglah $g \circ f$ dalam bentuk $z \mapsto az + b$ lalu
        kenalilah ia (jenis, pusat, sudutnya).
  \item Misalkan $r_1(z) = \conj z$ (yaitu pencerminan pada sumbu real) dan
        $r_2(z) = \iu\conj z$ (yaitu pencerminan pada garis $y = x$).
        Hitunglah $r_2 \circ r_1$ lalu periksalah pertanyaan 9 pada contoh
        ini.
\end{enumerate}

\medskip
\noindent\textbf{Bagian IV --- \omterm{def:b1:structures:group}{Grup} hingga: teorema Leonardo.}
Misalkan $G$ sebuah \omterm{def:b1:structures:group}{grup} \emph{hingga} berisi \omterm{def:b1:euclid:isometry}{isometri} bidang.
\begin{enumerate}[resume]
  \item Tunjukkan bahwa \omterm{def:b1:euclid:isometry}{isometri} mengawetkan barisentrum: bahwa jika
        $\lambda_1 + \dots + \lambda_m = 1$ dan $f = t_a \circ
        g$ (dengan $g$ \omterm{def:b1:linmaps:def}{linear}), maka $f\bigl(\sum_i \lambda_i
        x_i\bigr) = \sum_i \lambda_i f(x_i)$.
  \item Tetapkanlah $c = \frac{1}{\abs G}\sum_{g \in G} g(x_0)$ bagi sebarang
        $x_0$ yang dipilih. Tunjukkan bahwa setiap $h \in G$ menetapkan $c$:
        jadi \omterm{def:b1:structures:group}{grup} hingga berisi \omterm{def:b1:euclid:isometry}{isometri} mempunyai titik tetap bersama.
  \item Simpulkanlah bahwa, setelah dikonjugasikan oleh $t_{-c}$, orang boleh
        mengandaikan $G \subseteq O(\R^2)$. Misalkan $G^{+} = \{g \in G :
        \det g = 1\}$; tunjukkanlah bahwa entah $G = G^{+}$ atau $G^{+}$
        berindeks tepat $2$ di dalam $G$ (tunjukkanlah sebuah bijeksi $G^+
        \to G \setminus G^+$).
  \item Tunjukkan bahwa \omterm{def:b1:structures:group}{grup} hingga berisi \emph{rotasi} terhadap $c$
        bersifat siklik: di antara unsurnya pilihlah rotasi
        $R_{\theta_0}$ yang bersudut terkecil $\theta_0 \in
        \intoo{0}{2\pi}$, lalu pakailah pembagian Euclid atas sudutnya
        untuk membuktikan bahwa ia membangun; lalu simpulkanlah $\theta_0 =
        \frac{2\pi}n$ dan $G^{+} = \{R_{\theta_0}^{\,k}\} \cong
        C_n$.
  \item Andaikan $G \neq G^{+}$ lalu pilihlah sebuah pencerminan $s \in G$.
        Tunjukkan $G = G^{+} \cup sG^{+}$, bahwa $s r s = r^{-1}$ bagi
        setiap rotasi $r \in G^{+}$, dan bahwa kesemua $n$ unsur
        $sG^{+}$ merupakan pencerminan: jadi $G$ adalah \emph{\omterm{pb:b1:euclid:1}{grup
        dihedral}} $D_n$ berorde $2n$.
  \item Simpulkanlah (\emph{teorema Leonardo}): bahwa setiap \omterm{def:b1:structures:group}{grup} hingga
        berisi \omterm{def:b1:euclid:isometry}{isometri} bidang bersifat siklik $C_n$ atau dihedral
        $D_n$.
\end{enumerate}

\medskip
\noindent\textbf{Bagian V --- Dividen, dan sintesisnya.}
\begin{enumerate}[resume]
  \item Tunjukkan secara langsung bahwa \omterm{def:b1:structures:group}{grup} hingga berisi \omterm{def:b1:euclid:isometry}{isometri} tak dapat
        memuat translasi ataupun \omterm{pb:b1:euclid:1}{pencerminan geser} selain
        identitasnya (pandanglah pangkat unsur yang
        demikian).
  \item Misalkan $P_n$ segi $n$ beraturan yang bertitik sudut
        \omterm{def:b1:complex:unity}{akar satuan} ke-$n$ ($n \geq 3$). Tunjukkan bahwa \omterm{def:b1:structures:group}{grup}
        kesetangkupannya tepat $D_n$: yaitu $n$ rotasi $z
        \mapsto \omega^k z$ dan $n$ pencerminan $z \mapsto
        \omega^k \conj z$, dengan $\omega = \eu^{2\iu\pi/n}$, dan tak ada
        yang lain.
  \item Tunjukkanlah bangun bidang yang \omterm{def:b1:structures:group}{grup} kesetangkupannya $C_1$,
        $D_1$, $D_2$ dan $C_3$ berturut-turut.
  \item Daftarkanlah kedelapan unsur \omterm{def:b1:structures:group}{grup} kesetangkupan pada
        persegi yang bertitik sudut $\pm1, \pm\iu$ sebagai \omterm{def:b1:logic:map}{pemetaan} $z \mapsto
        \omega^k z$ atau $z \mapsto \omega^k\conj z$, lalu berikanlah
        sumbu masing-masing dari keempat pencerminannya.
  \item Misalkan $f, g$ rotasi bersudut sama $\theta
        \notin 2\pi\Z$ terhadap pusat yang \emph{berbeda} $c_1 \neq
        c_2$. Hitunglah $f\circ g - g\circ f$ titik demi titik lalu tunjukkanlah
        $f\circ g \neq g\circ f$; tunjukkanlah lebih-lebih bahwa $(f \circ
        g)\circ(g\circ f)^{-1}$ merupakan translasi yang tak sepele, sehingga
        sebarang \omterm{def:b1:structures:group}{grup} yang memuat $f$ dan $g$ takhingga ---
        yaitu penjelasan kedua atas pusat tunggal pada teorema
        Leonardo.
  \item Sintesis, dalam empat kalimat: kedua hasil struktur yang mana
        yang mereduksi \omterm{def:b1:euclid:isometry}{isometri} sebarang menjadi aljabar \omterm{def:b1:linmaps:def}{linear}
        (pertanyaan 3--4) dan \omterm{def:b1:structures:group}{grup} hingga sebarang menjadi
        \omterm{def:b1:structures:subgroup}{subgrup} $O(2)$ (pertanyaan 15); apakah daftar lengkap
        \omterm{def:b1:euclid:isometry}{isometri} bidangnya dan invarian yang mana
        (langsung/taklangsung, titik tetap) yang memisahkan keempat jenisnya;
        mengapakah kaidah komposisi pertanyaan 9 menjadikan pencerminan
        pembangun segalanya; dan apa yang ditambahkan teorema
        Leonardo pada skala hingganya. Namailah kedua teorema
        yang dibuktikan pada Bagian II dan IV.

\end{enumerate}
\end{problem}
